\section*{Chapitre \ref{ch:b1:crystals-metals} --- Cristaux I~: le cristal parfait et les métaux}

\begin{solution}{exo:b1:crystals-metals:1}
Cubique simple~: $8 \times \frac18 = 1$ atome, \omterm{def:b1:crystals-metals:coordination}{coordinence} 6. CC~: $8 \times
\frac18 + 1 = 2$, \omterm{def:b1:crystals-metals:coordination}{coordinence} 8. CFC~: $8 \times \frac18 + 6 \times
\frac12 = 4$, \omterm{def:b1:crystals-metals:coordination}{coordinence} 12.
\end{solution}

\begin{solution}{exo:b1:crystals-metals:2}
Les atomes sont tangents le long d'une arête~: $a = 2R$, $Z = 1$, $C = \frac43\pi R^3/(2R)^3
= \pi/6 \approx 0.52$, moins que le CC (0.68) et le CFC (0.74).
\end{solution}

\begin{solution}{exo:b1:crystals-metals:3}
$R = a\sqrt2/4 = 405.0 \times 0.3536 = \qty{143.2}{pm}$;
$\rho = 4 \times 26.98/(6.022 \times 10^{23} \times (4.050 \times
10^{-8})^3) = \qty{2.70}{g/cm^3}$.
\end{solution}

\begin{solution}{exo:b1:crystals-metals:4}
$R = a\sqrt3/4 = \qty{185.8}{pm}$; $\rho = 2 \times 22.99/(6.022 \times
10^{23} \times (4.291 \times 10^{-8})^3) = \qty{0.966}{g/cm^3}$, en
accord avec la valeur mesurée, \qty{0.97}{g/cm^3}.
\end{solution}

\begin{solution}{exo:b1:crystals-metals:5}
$a^3 = 4M/(\rho N_A) = 4 \times 107.87/(10.50 \times 6.022 \times 10^{23}) =
\qty{6.824e-23}{cm^3}$, donc $a = \qty{4.086e-8}{cm} = \qty{408.6}{pm}$ et
$R = a\sqrt2/4 = \qty{144.5}{pm}$.
\end{solution}

\begin{solution}{exo:b1:crystals-metals:6}
La démonstration est celle de la \cref{prop:b1:crystals-metals:hcp}.
Magnésium~: $c/a = 521.0/320.9 = 1.624$. Le prisme de volume
$\frac{3\sqrt3}{2}a^2c$ contient 6 atomes, donc $\rho = 6M/(N_A \cdot
\frac{3\sqrt3}{2}a^2c) = 6 \times 24.31/(6.022 \times 10^{23} \times
2.598 \times (3.209 \times 10^{-8})^2 \times 5.210 \times 10^{-8}) =
\qty{1.74}{g/cm^3}$.
\end{solution}

\begin{solution}{exo:b1:crystals-metals:7}
Atomes~: $(0,0,0)$, $(\frac a2,\frac a2,0)$, $(\frac a2,0,\frac a2)$,
$(0,\frac a2,\frac a2)$. \omterm{def:b1:crystals-metals:site}{Sites octaédriques}~: $(\frac a2,\frac a2,\frac a2)$,
$(\frac a2,0,0)$, $(0,\frac a2,0)$, $(0,0,\frac a2)$. \omterm{def:b1:crystals-metals:site}{Sites tétraédriques}~:
les 8 points dont chaque coordonnée vaut $\frac a4$ ou $\frac{3a}{4}$. Un
\omterm{def:b1:crystals-metals:site}{site octaédrique} et deux \omterm{def:b1:crystals-metals:site}{sites tétraédriques} par atome.
\end{solution}

\begin{solution}{exo:b1:crystals-metals:8}
$R(\ce{Cu}) = 361.5\sqrt2/4 = \qty{127.8}{pm}$, $R(\ce{Ni}) = 352.4\sqrt2/4
= \qty{124.6}{pm}$~: des rayons à 3\,\% près et la même structure, de
sorte que chaque atome peut prendre la place de l'autre sur le \omterm{def:b1:crystals-metals:lattice}{réseau}.
L'hydrogène, bien plus petit, ne peut se loger que dans les \omterm{def:b1:crystals-metals:site}{sites
interstitiels}~: il forme un \omterm{def:b1:crystals-metals:alloy}{alliage d'insertion}.
\end{solution}

\begin{solution}{exo:b1:crystals-metals:9}
$R = 316.5\sqrt3/4 = \qty{137.0}{pm}$~; $\rho = 2 \times 183.84/(6.022 \times
10^{23} \times (3.165 \times 10^{-8})^3) = \qty{19.26}{g/cm^3}$~; atomes
par \unit{cm^3}~: $2/a^3 = 2/(3.165 \times 10^{-8})^3 = \num{6.31e22}$.
\end{solution}

\begin{solution}{exo:b1:crystals-metals:10}
Le centre de face $(\frac a2,\frac a2,0)$ est à $a/2$ des centres des
\omterm{def:b1:crystals-metals:lattice}{mailles} situées au-dessus et au-dessous~: $r_O = \frac a2 - R = \frac a2 - \frac{a\sqrt3}{4}
= \frac{2-\sqrt3}{4}a$. Le point $(\frac a2,\frac a4,0)$ est à
$\sqrt{a^2/4 + a^2/16} = \frac{a\sqrt5}{4}$ des sommets $(0,0,0)$ et
$(a,0,0)$ et des centres $(\frac a2,\frac a2,\pm\frac a2)$~:
$r_T = \frac{\sqrt5 - \sqrt3}{4}a$. Avec $a = 4R/\sqrt3$~:
$r_O = \frac{2-\sqrt3}{\sqrt3}R = 0.155\,R$ et
$r_T = \frac{\sqrt5-\sqrt3}{\sqrt3}R = 0.291\,R$. Dans le CC, c'est le
\omterm{def:b1:crystals-metals:site}{site tétraédrique} le plus grand.
\end{solution}

\begin{solution}{exo:b1:crystals-metals:11}
$C = Z\cdot\frac43\pi R^3/a^3$ et $Z/a^3 = \rho N_A/M$, donc $C =
\frac{4\pi R^3\rho N_A}{3M} = \frac{\pi\rho N_A}{6M}(2R)^3$. Cuivre~:
$\pi \times 8.93 \times 6.022 \times 10^{23}/(6 \times 63.55) \times (2.556
\times 10^{-8})^3 = 0.740$.
\end{solution}

\begin{solution}{exo:b1:crystals-metals:12}
$a \approx (407.8 + 408.6)/2 = \qty{408.2}{pm}$~; masse molaire moyenne $(196.97 +
107.87)/2 = \qty{152.42}{g/mol}$~; $\rho = 4 \times 152.42/(6.022 \times
10^{23} \times (4.082 \times 10^{-8})^3) = \qty{14.9}{g/cm^3}$.
\end{solution}

\begin{solution}{pb:b1:crystals-metals:1}
\textbf{1.} $Z = 2$~; \omterm{def:b1:crystals-metals:coordination}{coordinence} 8.
\textbf{2.} Le long de la grande diagonale~: $a\sqrt3 = 4R$, $R = 286.65 \times
\sqrt3/4 = \qty{124.1}{pm}$.
\textbf{3.} $\rho = 2 \times 55.845/(6.022 \times 10^{23} \times (2.8665
\times 10^{-8})^3) = \qty{7.87}{g/cm^3}$.
\textbf{4.} $2/a^3 = 2/(2.8665 \times 10^{-8})^3 = \num{8.49e22}$ atomes.
\textbf{5.} $C = \pi\sqrt3/8 = 0.680$.
\textbf{6.} Les masses volumiques calculée et mesurée concordent à trois
chiffres.
\textbf{7.} $\alpha$~: $R = 289.8\sqrt3/4 = \qty{125.5}{pm}$~; $\gamma$~:
$R = 363.9\sqrt2/4 = \qty{128.7}{pm}$.
\textbf{8.} $\alpha$~: $a^3/2 = \qty{12.17e6}{pm^3}$ par atome~; $\gamma$~:
$a^3/4 = \qty{12.05e6}{pm^3}$ par atome.
\textbf{9.} Le fer $\gamma$ est plus dense~: $\rho_\alpha = 55.845/(6.022 \times 10^{23}
\times 1.217 \times 10^{-23}) = \qty{7.62}{g/cm^3}$, $\rho_\gamma =
\qty{7.70}{g/cm^3}$.
\textbf{10.} $(12.05 - 12.17)/12.17 = -1.0\,\%$~: la barre se contracte
d'environ 1\,\% quand on la chauffe au travers de la transition.
\textbf{11.} Oui~: le CFC est plus compact (0.74 contre 0.68), ce qui
l'emporte sur le rayon atomique plus grand~:
$(0.680/0.740) \times (128.7/125.5)^3 = 0.99$.
\textbf{12.} $r_C = a\sqrt3/8 = 356.68 \times 0.2165 = \qty{77.2}{pm}$.
\textbf{13.} $r_O = 0.0670 \times 289.8 = \qty{19.4}{pm}$.
\textbf{14.} $r_T = 0.1260 \times 289.8 = \qty{36.5}{pm}$.
\textbf{15.} Le carbone (\qty{77}{pm}) est deux fois plus gros que le
plus grand site~: les atomes de fer voisins doivent être écartés.
\textbf{16.} La déformation coûte beaucoup d'énergie~: très peu d'atomes
de carbone entrent dans le fer $\alpha$.
\textbf{17.} 4 \omterm{def:b1:crystals-metals:site}{sites octaédriques} par \omterm{def:b1:crystals-metals:lattice}{maille} de 4 atomes~: un par atome
de fer.
\textbf{18.} $r_T = (\sqrt{3/2} - 1) \times 128.7 = \qty{28.9}{pm}$.
\textbf{19.} Un carbone par fer~: \ce{FeC}, $w(\ce{C}) = 12.011/(55.845 +
12.011) = 17.7\,\%$.
\textbf{20.} Un carbone pour dix fers~: $w = 12.011/(558.45 + 12.011) =
2.1\,\%$.
\textbf{21.} Les \omterm{def:b1:crystals-metals:site}{sites octaédriques} du fer $\gamma$ sont bien plus grands
que tous les sites du fer $\alpha$~: le fer chaud dissout le carbone, que
la trempe piège ensuite.
\textbf{22.} $r_O = (\sqrt2 - 1) \times 128.7 = \textbf{\qty{53}{pm}}$, le
plus grand vide du fer, encore plus petit qu'un atome de carbone
(\qty{77}{pm})~: le carbone se dissout dans le fer chaud, mais en le
distendant.
\end{solution}
